\section{TapeAgents：可审计的执行轨迹}
Chapados 在 00:18:56--00:20:52 介绍 TapeAgents，这是一个刚刚开源的框架，目标是“在工程工具和优化工具之间搭建桥梁”。“Tape”是统一的抽象——所有 Agent 的思考、行动、观测都写进这条日志，既可审计，也可被后续 Agent 读取。

\begin{importantbox}{TapeAgents 的“Best of Both Worlds”}
``我们想要开发/调试工具，又想要具备 prompt 优化与 fine-tuning 的能力''；TapeAgents 提供细粒度组件、并行流、runtime orchestration，同时以 tape 作为数据资产，将执行日志转化为优化材料。
\end{importantbox}

\subsection{Tape 的结构与可重用性}
TapeAgents 里的 tape 不是纯粹的日志，它带有丰富的结构化 metadata，可以被其他 Agent 或优化器消费（00:21:10--00:22:25）。多个 Agent 可以共享同一 tape，Agent B 读取全历史并选择想要执行的 node，从而构建模块化的控制流。

\begin{knowledgebox}{Tape 作为“数据即服务”}
Tape 把每一次 prompt、每一次环境反馈、行动指令都记录下来，进而可以驱动调试、可视化、prompt 优化、教师-学生蒸馏，形成一条从 execution trace 到 fine-tune 数据的闭环。
\end{knowledgebox}

\subsection{本章小结}
TapeAgents 借助丰富 metadata 的统一 tape，实现了多 Agent 协作、工程级调试与自动化优化的协同，解决了“人工 prompt + 实验优化”之间的脱节。

